\documentclass{article}
\usepackage[letterpaper,top=2cm,bottom=2cm,left=3cm,right=3cm,marginparwidth=1.75cm]{geometry}

\usepackage{graphicx} % Required for inserting images
\usepackage{booktabs}
\usepackage{caption}
\usepackage{subcaption}
\usepackage{authblk}
\usepackage{amsmath}
\usepackage{amssymb}
\usepackage[dvipsnames]{xcolor}
\usepackage[english]{babel}
\usepackage{pgfgantt}
\usepackage{csquotes}
\usepackage{multirow}

\usepackage[colorlinks=true, allcolors=blue]{hyperref}
\hypersetup{colorlinks=true, breaklinks=true}

\usepackage[    
    style=numeric-comp,
    sorting=none,
    sortcites=true
    ]{biblatex}

\bibliography{latex/bibliography}

\newcommand{\mycitethis}{\tkDE{[?]}}

\usepackage{lscape}

\usepackage{verbatim}
\usepackage{geometry}


% Change font
\usepackage{helvet}
\renewcommand{\familydefault}{\sfdefault}

\usepackage{xcolor}
\newcommand{\tkDE}[1]{\textcolor{magenta}{#1}} 
\newcommand{\citethis}{\tkDE{[?]}}

\newcommand{\diff}{\mathrm{d}}  

\title{BSol Measles Participatory System Modelling (July Update)}
\author{David Ellis}
\date{30/04/2026}

\begin{document}

\maketitle

\section{Workshop Summary}

A a participatory measles modelling workshop was held on the 20th of March 2026. This workshop, hosted by Birmingham City Council, brought together local experts to build a conceptual model of the local measles system, to identify behavioural and epidemiological factors that facilitate and block flows in the system, and to identify potential interventions to improve these flows.

Due to the length of the workshop, Copilot was used to extract vocal contributions to supplement the written materials. The full output, along with all non-sensitive materials related to this project, can be viewed on  \href{https://github.com/David-Ellis/Measles-System-Dynamics-BSol}{GitHub}.

\tkDE{Brief text summary.}

The resulting models have been digitally transcribed and are presented in Appendix~\ref{app:trans-models} as Figures~\ref{fig:group1-model} and \ref{fig:group2-model}. Photographs of the models produced in the workshop can also be found on \href{https://github.com/David-Ellis/Measles-System-Dynamics-BSol}{GitHub}.

\begin{figure}
    \centering
    \includegraphics[width=0.8\linewidth]{figures/workshop/flow-factors.pdf}
    \caption{Caption}
    \label{fig:flow-factors}
\end{figure}

\section{Basic Modelling}

\begin{table}[ht!]
    \begin{tabular}{p{0.5\linewidth} | p{0.25\linewidth} | p{0.15\linewidth}}
    \toprule
    \textbf{Variable} & \textbf{Value} & \textbf{Source(s)}\\ \midrule
    
    Initial number of babies ($<$1 yrs) & 16,346 & \cite{ONS2021census} \\
    
    Initial number of children vaccinated (1 to 4 yrs) & 4,745&  \cite{ONS2021census, NHSEngland2026FDP}  \\
    
    Initial number of children unvaccinated (1 to 4 yrs) & 14,620 & \cite{ONS2021census, NHSEngland2026FDP} \\
    
    Birth rate & 16,346 yr$^{-1}$ & \cite{ONS2021census} \\
    
    Baby ($<$1 yrs) migration rate & 196 yr$^{-1}$ & \cite{ONS2021census} \\
    
    Vaccinated child (1 to 4 yrs) migration rate & 801.7 yr$^{-1}$ & \cite{ONS2021census, WHO_measles_coverage} \\
    
    Unvaccinated child (1 to 4 yrs) migration rate & 80.1 yr$^{-1}$ & \cite{ONS2021census, WHO_measles_coverage} \\

    MMR/MMRV vaccination rate (1 to 4 yrs) & 15,360 yr$^{-1}$ & \cite{NHSEngland2026FDP} \\

    Measles basic reproduction number, $R_0$ & 13.5 & \cite{https://doi.org/10.1111/risa.12637} \\

    Measles infectious period, $1/\gamma$ & 1.25 weeks & \cite{WHO2025MeaslesFactsheet} \\

    Measles transmission rate, $\beta = R_0\gamma$ & 1.54 day$^{-1}$ &  \cite{WHO2025MeaslesFactsheet, https://doi.org/10.1111/risa.12637}  \\

    Hospitalisation rate & 0.0021 day$^{-1}$ & \cite{NICE2024MeaslesComplications} \\

    Median hospital stay & 2 days & Locally calculated \\
    
    \end{tabular}
    \caption{Model parameter values for a basic measles system model (No ethnicity or behavioural factors).}
    \label{tab:modelling-papers} 
\end{table}

\subsection{SIR model}

\tkDE{Talk about initial SIR modelling and how the behaviour didn't match observations.}

\subsection{SIR + Isolation}

\tkDE{Had to introduce mechanism to prevent continuous outbreak cycle.}



\begin{figure}
    \centering
    \includegraphics[width=1\linewidth]{figures/model-diagrams/SIR+isolation.png}
    \caption{Diagram of a basic SIR system with the ability to isolate infected individuals at a fixed daily proportion once the number of infections reaches a fixed threshold for action.}
    \label{fig:SIR+isol}
\end{figure}

\begin{figure}
    \centering
    \includegraphics[width=0.9\linewidth]{figures/model-params/SIR-isol-thresh.pdf}
    \caption{Number of people infected each day for different isolation thresholds with $S(0)/N = 0.1$, $R_0 = 13.5$, and an isolation proportion of 0.1.}
    \label{fig:isol-thresh}
\end{figure}

\subsection{Full model}

\subsection{Model tuning}

To maximise the use of data, the model is tuned by minimising an error function. Since relative timing of the simulated is not \tkDE{important to this study}, a time-lag parameter $\tau$ is used so that the modelled outcomes are compared to observed outcomes occurring at the time time. The error function is then given by
%
\begin{equation}
    E_V(p_i, \tau) = \sum_t\left( V_{t+\tau}^{\text{sim}}(p_i) - V_t^{\text{obs}}\right)^2 
\end{equation}
%
where $p_i$ represents all free model parameters to be tuned, and $V_{t+\tau}^{\text{sim}}$ and $V^{\text{obs}}$ are the simulated and observed measurable values respectively. 

There are two observable values for the model to be tuned to: measles infections and hospital admissions. Therefore, the full error function to be minimised is
%
\begin{equation}
    E(p_i) = E_I(p_i) + \omega E_H(p_i)
\end{equation}
%
where $E_I$ and $E_H$ are the error functions for infections and hospital admissions respectively and $\omega$ is a weighting parameter that encodes the relative importance of $E_H$ to $E_I$ in the error function calculation. 

\printbibliography

\appendix

\newpage
\section{Transcribed Models}
\label{app:trans-models}

\begin{figure}[ht!]
    \centering
    \includegraphics[width=1\linewidth]{workshop/workshop-model-transcription-group1.pdf}
    \caption{Caption}
    \label{fig:group1-model}
\end{figure}

\begin{figure}[ht!]
    \centering
    \includegraphics[width=1\linewidth]{workshop/workshop-model-transcription-group2.pdf}
    \caption{Caption}
    \label{fig:group2-model}
\end{figure}

\end{document}